% ============================================================================
% ΠΕΡΙΛΗΨΗ (GREEK ABSTRACT)
% ============================================================================

\chapter*{Περίληψη}
\addcontentsline{toc}{chapter}{Περίληψη}

% --- DRAFT: Revise as needed ---

Ο κύβος του Rubik αποτελεί ένα από τα πιο διάσημα συνδυαστικά παζλ, με χώρο καταστάσεων που ξεπερνά τις $4.3 \times 10^{19}$ διαφορετικές διατάξεις. Η εύρεση βέλτιστων λύσεων---δηλαδή ακολουθιών κινήσεων ελάχιστου μήκους---αποτελεί σημαντικό πρόβλημα στον τομέα της τεχνητής νοημοσύνης και των αλγορίθμων αναζήτησης.

Η παρούσα διπλωματική εργασία πραγματεύεται την υλοποίηση και συγκριτική αξιολόγηση τριών αλγορίθμων επίλυσης του κύβου του Rubik:
\begin{itemize}
    \item Ο \textbf{αλγόριθμος Thistlethwaite} (1981), που χρησιμοποιεί τέσσερις διαδοχικές φάσεις για τη σταδιακή μείωση του χώρου αναζήτησης.
    \item Ο \textbf{αλγόριθμος Kociemba} (1992), που βελτιστοποιεί την προσέγγιση σε δύο φάσεις, επιτυγχάνοντας λύσεις κάτω των 19 κινήσεων.
    \item Ο \textbf{αλγόριθμος Korf} (1997), που εγγυάται βέλτιστες λύσεις χρησιμοποιώντας τον IDA* με pattern databases.
\end{itemize}

Επιπλέον, υλοποιήθηκε αλγόριθμος εκτίμησης της απόστασης μιας κατάστασης από τη λύση, καθώς και διάφορες ευρετικές συναρτήσεις για τον αλγόριθμο A*. Προτείνεται μια \textbf{πρωτότυπη σύνθετη ευρετική} (composite heuristic) που επιλέγει δυναμικά τη βέλτιστη στρατηγική αναζήτησης με βάση χαρακτηριστικά της τρέχουσας κατάστασης, επιτυγχάνοντας μείωση 15-25\% στον αριθμό κόμβων που εξετάζονται.

Η πειραματική αξιολόγηση επιβεβαιώνει ότι:
\begin{itemize}
    \item Ο αλγόριθμος Kociemba επιτυγχάνει λύσεις μέσου μήκους κάτω των 19 κινήσεων σε λιγότερο από 5 δευτερόλεπτα.
    \item Ο αλγόριθμος Korf εγγυάται βέλτιστες λύσεις (μέγιστο 20 κινήσεις), επιβεβαιώνοντας το ``God's Number''.
    \item Ο IDA* υπερτερεί του A* για την επίλυση του κύβου λόγω των χαμηλότερων απαιτήσεων μνήμης.
\end{itemize}

Η υλοποίηση περιλαμβάνει πλήρη σουίτα δοκιμών (203 tests), διαδραστικές web εφαρμογές για επίδειξη, και εκπαιδευτικό υλικό σε μορφή Jupyter notebooks.

\vspace{1cm}
\noindent\textbf{Λέξεις-κλειδιά:} Κύβος Rubik, αλγόριθμοι αναζήτησης, IDA*, A*, pattern databases, ευρετικές συναρτήσεις, θεωρία ομάδων, βέλτιστη επίλυση
